% Compile with: xelatex french-english-vowels.tex
% Alternatively: lualatex french-english-vowels.tex (full LuaLaTeX install).
% Unicode IPA glyphs are supplied by DejaVu Sans (free and widely available).
% This is a subset of the standard IPA vowel-symbol chart, not a plot of
% measured French/English vowel realizations. Positions are the IPA reference
% positions; broad transcription does not specify exact accent-specific quality.
% Scope: common French and traditional British/General American English
% transcription symbols. French /ɑ/ and /œ̃/ are retained for accents with
% those distinctions. Length and English diphthongs are not plotted.
% Reference: https://www.internationalphoneticassociation.org/content/ipa-vowels
% French: Fougeron & Smith (1993), Illustrations of the IPA: French.
\documentclass[tikz,border=8pt]{standalone}
\usepackage{fontspec}
\setsansfont{DejaVu Sans}
\renewcommand{\familydefault}{\sfdefault}
\usepackage{tikz}
\definecolor{ipablue}{HTML}{2D63C8}
\definecolor{ipagreen}{HTML}{036264}
\definecolor{ipared}{HTML}{9e0808}
\definecolor{ipapurple}{HTML}{51089e}
\definecolor{guidegray}{HTML}{7C8795}
\definecolor{labelgray}{HTML}{354052}

% Place an IPA pair around a reference point. Either symbol may be empty.
% Left = unrounded; right = rounded, including singleton cardinal vowels.
% \newcommand{\vowelpair}[4]{%
%   \node[point] at (#1,#2) {};
%   \if\relax\detokenize{#3}\relax\else
%     \node[vowel,anchor=east] at (#1-0.13,#2) {#3};
%   \fi
%   \if\relax\detokenize{#4}\relax\else
%     \node[vowel,anchor=west] at (#1+0.13,#2) {#4};
%   \fi
% }

\newcommand{\vowelpair}[6]{%
  \begingroup
  \def\firststyle{vowel}%
  \def\secondstyle{vowel}%
  \if\relax\detokenize{#3}\relax\else
    \edef\firststyle{#3-vowel}%
  \fi
  \if\relax\detokenize{#5}\relax\else
    \edef\secondstyle{#5-vowel}%
  \fi
  \node[point] at (#1,#2) {};
  \if\relax\detokenize{#4}\relax\else
    \node[\firststyle,anchor=east] at (#1-0.13,#2) {#4};
  \fi
  \if\relax\detokenize{#6}\relax\else
    \node[\secondstyle,anchor=west] at (#1+0.13,#2) {#6};
  \fi
  \endgroup
}

\begin{document}
\begin{tikzpicture}[
  x=1cm,y=0.82cm,
  fr-vowel/.style={font=\fontsize{25}{29}\selectfont,text=ipablue,
                fill=white,inner sep=2pt},
  en-vowel/.style={font=\fontsize{25}{29}\selectfont,text=ipared,
                fill=white,inner sep=2pt},
  vowel/.style={font=\fontsize{25}{29}\selectfont,text=ipapurple,
                fill=white,inner sep=2pt},
  point/.style={circle,fill=black,draw=white,line width=2pt,
                minimum size=7pt,inner sep=0pt},
  heading/.style={font=\large\bfseries,text=labelgray},
  rowlabel/.style={font=\normalsize,text=labelgray,anchor=east},
  note/.style={font=\footnotesize,text=labelgray,anchor=west},
  guide/.style={draw=guidegray,line width=0.8pt}
]
  % Standard vowel trapezium: front, central and back reference lines.
  \draw[guide] (0,6) -- (7.2,6) -- (7.2,0) -- (3.2,0) -- cycle;
  \draw[guide] (3.6,6) -- (5.2,0);
  \draw[guide] (1.0667,4) -- (7.2,4);
  \draw[guide] (2.1333,2) -- (7.2,2);

  \node[heading] at (0.35,6.85) {Front};
  \node[heading] at (3.6,6.85) {Central};
  \node[heading] at (7.2,6.85) {Back};
  \foreach \y/\name in {
    6/Close,5/Near-close,4/Close-mid,3/Mid,
    2/Open-mid,1/Near-open,0/Open}
    \node[rowlabel] at (-0.85,\y) {\name};

  \vowelpair{0}{6}{}{i}{fr}{y}
  \vowelpair{7.2}{6}{}{}{}{u}
  \vowelpair{1.65}{5}{en}{ɪ}{}{}
  \vowelpair{6.05}{5}{}{}{en}{ʊ}
  \vowelpair{1.0667}{4}{fr}{e}{fr}{ø}
  \vowelpair{7.2}{4}{}{}{fr}{o};
  \node[vowel] at (4.4,3) {ə};
  \vowelpair{2.1333}{2}{}{ɛ}{fr}{œ}
  \vowelpair{7.2}{2}{en}{ʌ}{}{ɔ}
  \vowelpair{2.6667}{1}{en}{æ}{}{};
  \vowelpair{3.2}{0}{fr}{a}{}{};
  \vowelpair{7.2}{0}{}{ɑ}{en}{ɒ}

  \node[note] at (-2.8,-0.9)
  {At each point: left = unrounded; right = rounded.};
  \node[note] at (-2.8, -1.45)
  {Color: {\color{ipablue}French}, {\color{ipared}{English}}, {\color{ipapurple}both} };
  % \node[note] at (-2.8,-2)
  %   {American English r-coloured vowels: {\color{ipablue}ɚ\quad ɝ}.};
  % \node[note] at (-2.8,-2.55)
  %   {IPA reference positions. Inventories vary by accent; length and diphthongs omitted.};
\end{tikzpicture}
\end{document}
